\documentclass[pdflatex,sn-mathphys-num,iicol]{sn-jnl}
\usepackage[utf8]{inputenc}
\usepackage{graphicx}%
\usepackage{multirow}%
\usepackage{amsmath,amssymb,amsfonts}%
\usepackage{amsthm}%
\usepackage{mathrsfs}%
\usepackage[title]{appendix}%
\usepackage{xcolor}%
\usepackage{textcomp}%
\usepackage{manyfoot}%
\usepackage{booktabs}%
\usepackage{algorithm}%
\usepackage{algorithmicx}%
\usepackage{algpseudocode}%
\usepackage{listings}%
\usepackage{verbatim}%

\usepackage{float}
\usepackage{stfloats}
\theoremstyle{thmstyleone}%
\newtheorem{theorem}{Theorem}%
\newtheorem{proposition}[theorem]{Proposition}%

\theoremstyle{thmstyletwo}%
\newtheorem{example}{Example}%
\newtheorem{remark}{Remark}%

\theoremstyle{thmstylethree}%
\newtheorem{definition}{Definition}%

\raggedbottom


\begin{document}

\title[Article Title]{JamVote - Collaborative Music Room}

\author*[1]{\fnm{Deepam} \sur{Ahuja}}\email{deepam.mitmpl2023@learner.manipal.edu}

\affil*[1]{\orgdiv{School of Computer Engineering}, \orgname{Manipal Institute of Technology},
\orgname{Manipal Academy of Higher Education},
\orgaddress{\city{Manipal}, \postcode{576104}, \state{Karnataka}, \country{India}}}

\abstract{Collaborative music listening in group settings suffers from well-documented
governance failures: existing platforms permit arbitrary queue manipulation by individual
users, provide no structured collective feedback mechanism, and make no distinction between
active and inactive participants. JamVote is a native Android application that addresses
these limitations through a non-blocking democratic queue governance model built on Firebase
Firestore. The system assigns each queued song a dynamic score computed from upvotes,
downvotes, and the contributor's historical reputation, with the queue ordered
dynamically by vote score in real time across all connected clients. The host has the ability to skip the currently playing song at any time, immediately
advancing the queue to the next highest-voted song. A
session-end gamification layer computes per-user reputation deltas and distributes
badges across three award categories. The proposed system
demonstrates that lightweight distributed systems design, without a dedicated backend
server, can sustain real-time democratic queue management across concurrent mobile clients.}

\keywords{Android Application Development, Real-Time Systems, Dynamic Queue Management,
Non-Blocking Voting Systems, Music Recommendation Algorithms, Gamification,
Engagement-Based Ranking, Distributed Systems}

\maketitle

\newpage

%% ─────────────────────────────────────────────
%%  SECTION 1 — INTRODUCTION
%% ─────────────────────────────────────────────
\section{Introduction}\label{sec:intro}

\subsection{Background and Importance}

Shared music listening has long been a cornerstone of social interaction, from communal
gatherings to remote friend groups connected over the internet. The scale of this shift
toward digital listening is substantial: according to the International Federation of the
Phonographic Industry, paid music streaming subscription accounts reached 752 million
globally in 2024, growing at 10.6\% year-on-year, with streaming now accounting for
69\% of all recorded music revenues worldwide~\cite{ifpi_gmr_2025}. The proliferation of
smartphones and cloud-based streaming services has made collaborative listening
technically feasible; platforms such as Spotify, Apple Music, and YouTube Music offer
rudimentary shared playlist features. However, the social dynamics of group listening
introduce challenges that go beyond simple playlist sharing. In a group setting, musical
preference is inherently heterogeneous, and the governance of a shared queue directly
affects the experience of every participant. When one user can arbitrarily reorder or
dominate a queue, the collaborative experience degrades into single-user control with a
passive audience.

The demand for fairer, more engaging collaborative listening systems is evidenced by the
growing popularity of DJ-style party applications and crowd-sourced jukebox concepts.
Yet, existing solutions lack the infrastructure to dynamically balance democratic
participation, fairness constraints, playback continuity, and social engagement within a
single, lightweight mobile application.

\subsection{Problem Context and Practical Significance}

Existing collaborative music platforms exhibit several well-documented limitations. First,
they permit manual queue reordering, enabling any participant to override the collective
preference. Second, they provide no structured mechanism for collective feedback; users
cannot signal approval or disapproval of queued songs in real time. Third, when a shared
queue is exhausted, playback falls back to an individual user's library, breaking the
collaborative experience. Fourth, these systems make no distinction between active and
inactive participants, allowing a user who has left the session to retain influence over
the queue. Finally, there is no incentive structure to encourage high-quality music
curation; all contributions, regardless of reception, are treated equally.

These limitations result in sessions that are either chaotic due to unconstrained
manipulation, or stagnant due to the absence of automated continuity mechanisms. The
practical significance of addressing these gaps extends to parties, corporate events,
remote collaborative work sessions, and any environment where a group shares an audio
space.

\subsection{Research Gap}

A review of the existing literature reveals that prior work addresses isolated components
of this problem without integrating them into a cohesive system. Recommender systems for
groups~\cite{jamvote_group_context_mdpi, jamvote_recommendation_groups} propose
aggregation strategies but do not implement real-time queue control or session governance.
Voting-based systems such as Liquid FM~\cite{jamvote_arxiv_social_music} demonstrate
collective ranking but impose computational complexity that limits lightweight mobile
deployment. Studies on collaborative playlist
behaviour~\cite{jamvote_collaborative_playlists_study} confirm user preference for
lightweight governance but provide no algorithmic implementation. Gamification and
reputation research~\cite{jamvote_online_community_motivation, jamvote_gamification_works}
demonstrates that visible reward systems improve engagement but does not evaluate
manipulation resistance in competitive voting environments.

Critically, no existing work combines real-time democratic queue management,
engagement-weighted scoring with fairness constraints, active user tracking for adaptive
governance, and a gamification layer within a single deployable mobile application. This
constitutes the central gap that JamVote addresses.

\subsection{Research Question}

This work addresses the following research question: \textit{How can a mobile application
implement a non-blocking, democratic music queue system that dynamically balances
collective engagement, user fairness, and playback continuity without requiring unanimous
participation or a centralised moderator?}

\subsection{Research Objectives}

The objectives of this project are as follows:

\begin{enumerate}
    \item To design and implement a real-time collaborative music Android application
    using Firebase Firestore as the synchronisation backend, supporting both private and
    public listening rooms.

    \item To implement a non-blocking democratic queue management system in which
    playback continuity is guaranteed regardless of participant voting activity.

    \item To develop a dynamic song scoring mechanism that ranks queued songs based on
    upvotes, downvotes, and the reputation weight of the song's contributor, with queue
    reordering triggered only on song completion to prevent disruptive mid-session
    reshuffling.

    \item To implement active user tracking via periodic presence heartbeats, enabling
    the system to reflect only currently engaged participants in the active user count.

    \item To design and integrate a gamification and reputation system that incentivises
    high-quality music curation through session-end awards and persistent reputation
    scores.

    \item To evaluate the system through functional testing across multiple concurrent
    users on physical Android devices.
\end{enumerate}

\subsection{Research Contribution}

This work makes the following contributions:

\begin{enumerate}
    \item \textbf{Non-blocking democratic queue governance}: A voting-based queue
    management system in which playback is never stalled by insufficient votes. Songs
    advance automatically when their duration expires, and the host can skip the
    currently playing song at any time to advance the queue.

    \item \textbf{Real-time score-based queue ordering}: Queued songs are ordered
    dynamically by their computed vote score via a live Firestore query, ensuring that
    the highest-voted song is always displayed at the top and selected for playback
    next without any manual intervention.

    \item \textbf{Vote deduplication}: Each vote is stored as a Firestore document with
    a composite key of the form \texttt{\{userId\}\_\{songId\}}, enforcing the
    one-vote-per-user-per-song constraint at the database level without requiring a
    query, and preventing double voting or vote fabrication.

    \item \textbf{Active presence-aware governance}: A lightweight heartbeat-based
    presence tracking system that maintains a real-time active user count, reflecting
    only participants who are currently engaged in the session.

    \item \textbf{Full-stack gamification}: A session-end reputation system that
    rewards curators based on upvotes and downvotes received on their contributed
    songs, with a participation bonus applied for active contributors, persisted
    across sessions in a cloud database.
\end{enumerate}

\subsection{Paper Roadmap}

The remainder of this report is structured as follows.
Section~\ref{sec:litreview} reviews related work in collaborative playlist systems,
group recommender systems, and engagement-driven platforms, and identifies the gaps that
motivate this work. Section~\ref{sec:methodology} describes the system architecture,
data model, scoring algorithm, and implementation methodology.
Section~\ref{sec:results} presents the results of functional evaluation across
multi-device test scenarios. Section~\ref{sec:conclusion} summarises the findings and
restates the project contributions. Section~\ref{sec:limitations} discusses the
constraints and scope boundaries of this study. Section~\ref{sec:future} outlines
directions for future development.

%% ─────────────────────────────────────────────
%%  SECTION 2 — LITERATURE SURVEY
%% ─────────────────────────────────────────────
\section{Literature Survey}\label{sec:litreview}

The following subsections review work across five relevant areas: collaborative playlist
behaviour, automated recommendation, group recommender systems, prior collaborative music
applications, and gamification. Table~\ref{tab:litreview} provides a consolidated
comparative summary.

%% TABLE 1 — placed HERE, before subsections, so LaTeX floats it
%% within the Literature Survey section rather than after Methodology.
\begin{table*}[!ht]
\centering
\caption{Comparative Analysis of Related Work}
\label{tab:litreview}
\begin{tabular}{p{2.8cm} p{0.9cm} p{3.0cm} p{3.5cm} p{3.5cm}}
\toprule
\textbf{Author(s)} & \textbf{Year} & \textbf{Technique / Model} & \textbf{Key Contribution} & \textbf{Limitation} \\
\midrule

Park et al.~\cite{jamvote_collaborative_playlists_study}
& 2023
& Cross-cultural survey and statistical analysis
& Identifies user preference for lightweight governance in collaborative playlists
& No algorithmic implementation; inactive user differentiation not addressed \\

\addlinespace

Patel et al.~\cite{jamvote_integrated_playlist}
& 2023
& Hybrid collaborative and content-based filtering
& Automated playlist expansion with improved recommendation relevance
& Assumes single-user ownership; multi-user conflict resolution not addressed \\

\addlinespace

Chen~\cite{jamvote_cf_svd_ieee}
& 2023
& Collaborative filtering with SVD
& Improves recommendation accuracy via latent feature extraction
& Targets individual personalisation; synchronous group environments not handled \\

\addlinespace

Valera et al.~\cite{jamvote_group_context_mdpi}
& 2021
& Context-aware group aggregation
& Improves group satisfaction using fairness-aware aggregation strategies
& Evaluated on simulated data; real-time voting and playback continuity not implemented \\

\addlinespace

Liquid FM~\cite{jamvote_arxiv_social_music}
& 2015
& Voting-based viscous democracy model
& Collective vote-based music ranking via trust delegation
& Computational complexity limits mobile deployment; no playback continuity \\

\addlinespace

Jameson \& Smyth~\cite{jamvote_recommendation_groups}
& 2007
& Group aggregation (Borda count, least misery)
& Demonstrates that adaptive aggregation improves fairness perception
& Live session queue control and active user tracking not addressed \\

\addlinespace

Cunningham et al.~\cite{jamvote_playlist_creation_art}
& 2006
& Qualitative playlist analysis
& Identifies thematic and emotional logic in playlist construction
& No computational model proposed; conflict resolution not addressed \\

\addlinespace

Lampe et al.~\cite{jamvote_online_community_motivation}
& 2010
& Community participation analysis
& Demonstrates that reputation systems significantly improve sustained engagement
& Manipulation resistance in competitive voting contexts not evaluated \\

\addlinespace

Hamari et al.~\cite{jamvote_gamification_works}
& 2014
& Systematic review of 24 empirical gamification studies
& Establishes that gamification produces positive engagement effects in appropriate contexts
& Effectiveness highly context-dependent; adversarial voting settings not considered \\

\addlinespace

Mirea~\cite{jamvote_queue_mirea}
& 2019
& Host-controlled collaborative playlist mobile app (Spotify)
& Demonstrates feasibility of multi-user collaborative song addition on mobile
& No voting or scoring mechanism; no fairness constraints, presence tracking, or gamification \\

\bottomrule
\end{tabular}
\end{table*}

\subsection{Collaborative Playlist Behaviour}

Park et al.~\cite{jamvote_collaborative_playlists_study} found through cross-cultural
survey that users prefer lightweight governance mechanisms and that excessive control
reduces engagement, but provide no algorithmic implementation. Cunningham et
al.~\cite{jamvote_playlist_creation_art} established through qualitative study that
playlist construction follows thematic and emotional logic rather than purely algorithmic
optimisation, but propose no computational model for conflict resolution between
concurrent contributors.

\subsection{Automated Playlist Expansion and Individual Recommendation}

Patel et al.~\cite{jamvote_integrated_playlist} proposed a hybrid collaborative and
content-based filtering system for playlist expansion, improving recommendation relevance
but assuming single-user ownership with no mechanism for multi-user conflict resolution.
Chen~\cite{jamvote_cf_svd_ieee} applied SVD to a user-item interaction matrix for music
recommendation, achieving improved accuracy over neighbourhood-based methods but targeting
individual personalisation rather than synchronous group listening.

\subsection{Group Recommender Systems and Aggregation}

Valera et al.~\cite{jamvote_group_context_mdpi} evaluated context-aware aggregation
strategies including least misery and fairness-aware methods, demonstrating improved group
satisfaction on simulated datasets but without real-time voting or playback continuity.
Jameson and Smyth~\cite{jamvote_recommendation_groups} studied Borda count and weighted
averaging for group recommendation, showing that adaptive aggregation improves fairness
perception, but did not address live session queue control or active user tracking. The
Liquid FM model~\cite{jamvote_arxiv_social_music} introduced viscous democracy-based
music ranking via vote delegation, improving on majority voting but with computational
complexity that limits mobile deployment and no playback continuity mechanism.

\subsection{Prior Collaborative Music Applications}

Mirea~\cite{jamvote_queue_mirea} implemented a host-controlled collaborative playlist
mobile application connected to Spotify, demonstrating the feasibility of multi-user song
addition on mobile. However, the system provides no voting or scoring for queue
prioritisation, no fairness constraints, no active user tracking, and no gamification
layer.

\subsection{User Participation and Gamification}

Lampe et al.~\cite{jamvote_online_community_motivation} demonstrated through large-scale
community analysis that visible reputation systems significantly improve sustained
engagement, but did not evaluate manipulation resistance in competitive voting contexts.
Hamari, Koivisto, and Sarsa~\cite{jamvote_gamification_works} conducted a systematic
review of 24 empirical gamification studies, finding generally positive effects on
engagement but noting that effectiveness is highly context-dependent and moderated by the
quality of reward design.

\subsection{Research Gap}

The reviewed literature reveals two persistent and unresolved gaps. First, a
\textit{methodological gap}: existing group recommender
systems~\cite{jamvote_group_context_mdpi, jamvote_recommendation_groups} and voting
models~\cite{jamvote_arxiv_social_music} evaluate aggregation strategies on static
datasets without implementing real-time playback continuity or session-level governance,
making their findings inapplicable to live collaborative environments. Second, a
\textit{practical-knowledge gap}: the only directly comparable mobile
deployment~\cite{jamvote_queue_mirea} demonstrates multi-user song addition but omits
voting, scoring, fairness constraints, and gamification entirely, while gamification
research~\cite{jamvote_gamification_works, jamvote_online_community_motivation}
establishes engagement benefits without evaluating manipulation resistance in adversarial,
real-time voting contexts. Together, these gaps mean no existing system can simultaneously
guarantee playback continuity, enforce democratic fairness, adapt governance to active
participant count, and incentivise quality curation within a single lightweight mobile
application. JamVote addresses this by integrating all four capabilities into a single
serverless Android application built on Firebase Firestore, directly bridging the divide
between theoretical aggregation research and practical deployable governance.

%% ─────────────────────────────────────────────
%%  SECTION 3 — METHODOLOGY
%% ─────────────────────────────────────────────

\clearpage
\section{Methodology}\label{sec:methodology}

\subsection{System Overview}

JamVote is implemented as a native Android application using Java, XML-based layouts,
and a Groovy DSL Gradle build configuration, targeting Android API 25 (Android 7.1
Nougat) as the minimum SDK and API 34 as the target. The application uses Firebase
Authentication for identity management and Cloud Firestore as the sole backend data
store, eliminating the need for a dedicated application server. All client devices
communicate exclusively through Firestore snapshot listeners, which push document
updates to all connected clients within approximately one to two seconds of a write
operation. The host device additionally manages playback via a YouTube IFrame API
embedded in a \texttt{WebView}, with non-host clients synchronising to the currently
playing video by observing the \texttt{nowPlaying} field on the room document.

The overall system is composed of six functional modules: User Authentication and Room
Management, Real-Time Collaborative Queue, Non-Blocking Voting, Presence and Active
User Tracking, Playback Management, and Gamification and Reputation. These modules
interact exclusively through Firestore documents and are mediated by a centralised
\texttt{FirebaseUtil} class, which is the sole point of contact between application
logic and the Firebase SDK. No Activity or ViewModel calls the Firebase SDK directly.

\subsection{System Architecture}

Figure~\ref{fig:arch} illustrates the high-level architecture of JamVote. The client
layer consists of Android devices running the JamVote application. Each client maintains
a persistent Firestore snapshot listener for the room document, the queue
sub-collection, and the presence sub-collection. Write operations (votes, song
additions, presence heartbeats) are issued by individual clients through
\texttt{FirebaseUtil} and propagate to all other clients via Firestore's real-time
notification infrastructure.


The architecture follows a strict layering policy: Activities observe LiveData exposed
by ViewModels, ViewModels issue calls to \texttt{FirebaseUtil}, and
\texttt{FirebaseUtil} is the only class that interacts with the Firestore SDK. This
enforces separation of concerns and ensures that Firebase error handling is centralised.
All Firebase operations are asynchronous and use callback interfaces
(\texttt{AuthCallback}, \texttt{RoomCallback}, \texttt{SimpleCallback}, etc.) to
propagate results back to the ViewModel layer.

Music search is handled by the \texttt{YouTubeProvider} class, which issues search
queries to the YouTube Data API v3. A second batch API call to the \texttt{videos}
endpoint retrieves \texttt{contentDetails} for the returned videos and performs
explicit content detection by checking whether
\texttt{contentDetails.contentRating.ytRating} equals \texttt{ytAgeRestricted}.

\begin{figure}[!ht]
\centering
\includegraphics[width=0.85\linewidth]{arch_diag.png}
\caption{High-level system architecture. All clients communicate through Cloud Firestore;
no dedicated application server is required.}
\label{fig:arch}
\end{figure}

\subsection{Data Model}

The Firestore data model is organised into a top-level \texttt{users} collection and a
top-level \texttt{rooms} collection. The \texttt{rooms} collection contains three
sub-collections per room: \texttt{queue}, \texttt{votes}, and \texttt{presence}.
Table~\ref{tab:schema} summarises the key document structures.

%% TABLE 2 — placed here, inside Data Model subsection where it is referenced.
\begin{table*}[!ht]
\centering
\caption{Firestore document schema summary}
\label{tab:schema}
\begin{tabular}{p{2.0cm} p{2.2cm} p{4.8cm}}
\toprule
\textbf{Collection} & \textbf{Key Fields} & \textbf{Purpose} \\
\midrule
\texttt{users} & \texttt{uid}, \texttt{reputation}, \texttt{badges} &
Persistent user profile and cross-session reputation \\
\addlinespace
\texttt{rooms} & \texttt{roomCode}, \texttt{hostUid}, \texttt{status},
\texttt{sfwMode} &
Room state, host identity, content filter flag \\
\addlinespace
\texttt{queue} & \texttt{youtubeVideoId}, \texttt{score}, \texttt{position},
\texttt{addedByUid} &
Ordered song queue with computed scores \\
\addlinespace
\texttt{votes} & \texttt{userId}, \texttt{songId}, \texttt{value} &
One document per user-song pair; composite key prevents double voting \\
\bottomrule
\end{tabular}
\end{table*}

The \texttt{votes} sub-collection uses a composite document ID of the form
\texttt{\{userId\}\_\{songId\}}, which enforces the one-vote-per-user-per-song
constraint at the database level without requiring a query.

Active user presence is tracked using Firebase Realtime Database (RTDB) rather than
Firestore. Each client writes its display name to the path
\texttt{presence/\{roomId\}/\{uid\}} in RTDB on joining and on each heartbeat tick.
An \texttt{onDisconnect().removeValue()} handler is registered once per session,
causing the RTDB entry to be removed automatically when the client loses connectivity,
ensuring disconnected users are no longer counted as active. As a complementary
staleness filter, a user is classified as active only if their presence entry was
updated within the following window:

\begin{equation}
\text{lastSeen} > t_{\text{now}} - 90\text{ seconds}
\label{eq:presence}
\end{equation}

Each client writes a heartbeat every 30 seconds, ensuring that connected clients
consistently satisfy this condition while unresponsive clients are filtered out even
if their RTDB entry has not yet been removed.

\subsection{Core Algorithm Design}

\subsubsection{Dynamic Score Formula}

Each song in the queue is assigned a dynamic score used to determine playback order.
The score is defined as:

\begin{equation}
\text{score} = (\text{upvotes} - \text{downvotes}) +
\left\lfloor \frac{\text{adder\_reputation}}{10} \right\rfloor
\label{eq:score}
\end{equation}

where \texttt{adder\_reputation} is the persistent reputation score of the user who
added the song, not the users who voted. The floor division by 10 ensures that
reputation acts as a minor tiebreaker rather than a dominant factor, preserving the
primacy of collective voting. The score is recomputed and written to Firestore after
every vote event by the voting user's client.

\subsubsection{Dynamic Score-Based Queue Ordering}

The queue sub-collection is queried with an \texttt{orderBy("score", DESCENDING)}
clause, causing Firestore to re-deliver the ordered result set to all connected clients
whenever any song's score changes. The queue list therefore updates live on every vote
across all devices without additional client-side logic.

On song completion, the host client additionally performs a local
\texttt{Collections.sort()} over the remaining queued songs by score and writes
updated \texttt{position} values to Firestore in a single batch, maintaining the
\texttt{position} field as a persistent record of playback order.


\subsubsection{Host Skip}

The host has exclusive authority to skip the currently playing song. A skip button
is rendered in \texttt{QueueActivity} only for the host. Pressing it calls
\texttt{reorderAndPopNextSong()}, which marks the current song as played, sorts the
remaining queued songs by score, and begins playback of the top-ranked song.
Guest clients have no skip mechanism.

\subsection{Algorithm Pseudocode}

Algorithm~\ref{alg:songend} describes the song-end procedure executed by the host
client, covering the score-based sort and queue update.

\begin{algorithm}
\caption{Host Song-End Procedure}
\label{alg:songend}
\begin{algorithmic}[1]
\Require Current song $s_{\text{curr}}$, queue $Q$ as list of Song objects
\Ensure Firestore queue updated with new order; next song begins playback

\Function{OnSongEnded}{$s_{\text{curr}}$, $Q$}
    \State $Q \gets Q \setminus \{s_{\text{curr}}\}$
    \Comment{Remove completed song from queue}
    \State \Call{RemoveSongFromFirestore}{$s_{\text{curr}}$}
    \If{$Q = \emptyset$}
        \State Display empty queue notice; \Return
    \EndIf
    \State \Call{Sort}{$Q$, descending by \texttt{score}}
    \For{$i \gets 0$ \textbf{to} $|Q|-1$}
        \State $Q[i].\texttt{position} \gets i$
    \EndFor
    \State \Call{BatchWritePositions}{$Q$}
    \Comment{Single Firestore batch write}
    \State \Call{BeginPlayback}{$Q[0].\texttt{youtubeVideoId}$}
    \State \Call{WriteNowPlaying}{$Q[0].\texttt{youtubeVideoId}$}
    \Comment{Guests sync via listener}
\EndFunction
\end{algorithmic}
\end{algorithm}

Algorithm~\ref{alg:vote} describes the vote-casting procedure. Each user may cast
one vote per song; a second vote attempt on the same song is rejected by the
transaction.

\begin{algorithm}
\caption{Vote Casting Procedure}
\label{alg:vote}
\begin{algorithmic}[1]
\Require User $u$, song $s$, vote direction $v \in \{+1, -1\}$
\Ensure Firestore vote document and song score updated atomically

\Function{CastVote}{$u$, $s$, $v$}
    \State $\texttt{voteId} \gets u.\texttt{uid} + \text{``\_''} + s.\texttt{songId}$
    \State $\texttt{existing} \gets$ \Call{ReadVote}{\texttt{voteId}}
    \If{$\texttt{existing} = \emptyset$}
        \State \Call{WriteVote}{\texttt{voteId}, $v$}
        \State \Call{IncrementField}{$s$, $v = +1$ ? \texttt{upvotes} : \texttt{downvotes}}
        \State \Call{IncrementScore}{$s$, $v$}
        \Comment{Atomically update score}
    \Else
        \State \textbf{abort transaction}
        \Comment{One vote per user per song; duplicate rejected}
    \EndIf
\EndFunction
\end{algorithmic}
\end{algorithm}

\subsection{Navigation and Activity Structure}

The application navigation follows a linear host/guest bifurcation from a common entry
point. Figure~\ref{fig:nav} illustrates the complete navigation graph.

\begin{figure}[ht]
\centering
\fbox{\parbox{0.9\linewidth}{\small\centering
\texttt{SplashActivity} \\
$\downarrow$ \\
\texttt{LoginActivity} / \texttt{RegisterActivity} \\
$\downarrow$ \\
\texttt{MainActivity} (Create Room \textbar{} Join Room) \\[4pt]
$\swarrow$ \hspace{60pt} $\searrow$ \\
\texttt{CreateRoomActivity} \hspace{20pt} \texttt{JoinRoomActivity} \\
$\searrow$ \hspace{60pt} $\swarrow$ \\[4pt]
\texttt{QueueActivity} (host \textbar{} guest) \\
$\downarrow$ \hspace{30pt} $\downarrow$ \hspace{30pt} $\downarrow$ \\
\texttt{SearchActivity} \quad \texttt{LeaderboardActivity} \quad
\texttt{SessionSummaryActivity}
}}
\caption{Application navigation structure. All transitions use Android
\texttt{Intent}-based navigation.}
\label{fig:nav}
\end{figure}

\begin{figure*}[!ht]
\centering
\begin{tabular}{ccccc}
    \includegraphics[width=0.17\linewidth]{login.png} &
    \includegraphics[width=0.17\linewidth]{register.png} &
    \includegraphics[width=0.17\linewidth]{home_page.png} &
    \includegraphics[width=0.17\linewidth]{room_create.png} &
    \includegraphics[width=0.17\linewidth]{room_join.png} \\
    (a) Login & (b) Register & (c) Home & (d) Create room & (e) Join room \\
\end{tabular}
\caption{JamVote application screens: authentication flow and room management.}
\label{fig:ui_auth}
\end{figure*}

\begin{figure}[!ht]
\centering
\includegraphics[width=0.45\linewidth]{song_queue_list.png}
\caption{Collaborative queue screen showing dynamic song scoring, upvote/downvote
controls.}
\label{fig:ui_queue}
\end{figure}

\subsection{Setup}

The application was developed and tested under the following configuration:

\begin{itemize}
    \item \textbf{Development environment}: Android Studio Electric Eel (2022.1.1),
    Java 11, Groovy DSL \texttt{build.gradle}
    \item \textbf{Target platform}: Android API 34; minimum supported API 25
    (Android 7.1.1 Nougat)
    \item \textbf{Dependencies}: Firebase Authentication 22.3.1, Cloud Firestore
    24.10.3, Glide 4.16.0, AndroidX Lifecycle ViewModel and LiveData 2.7.0,
    Material Components 1.11.0
    \item \textbf{Music search}: YouTube Data API v3, accessed via HTTPS with the API
    key stored in \texttt{local.properties} and exposed through
    \texttt{BuildConfig.YOUTUBE\_API\_KEY}
    \item \textbf{Test devices}: Two physical Android devices running API 25 and API 33
    respectively, connected to the same Firebase project under a single Firestore
    database instance in the \texttt{asia-south1} region
    \item \textbf{Network}: Standard Wi-Fi; no network throttling applied during
    functional testing
\end{itemize}

\subsection{SFW Content Filtering}

The application supports a Safe-for-Work (SFW) mode toggled by the host at room
creation or mid-session. When SFW mode is active, a song is blocked from being added
to the queue if the song is flagged as explicit via \texttt{song.isExplicit()}. For the
YouTube provider, this flag is set by checking whether
\texttt{contentDetails.contentRating.ytRating} equals \texttt{ytAgeRestricted} in the YouTube
Data API \texttt{videos} endpoint response. The check is performed client-side before
the \texttt{addSong()} call reaches Firestore, and a \texttt{Snackbar} message is
displayed to the user. This is acknowledged as a best-effort filter, as YouTube's
content labelling is imperfect and unlabelled explicit content may pass through.

\subsection{Gamification and Reputation Computation}

At session end, the host triggers a status update setting \texttt{room.status =
"ended"}, which all clients observe via their room snapshot listener and use to
navigate to \texttt{SessionSummaryActivity}. The host computes the following per-user
statistics from the session's song history sub-collection:

\begin{itemize}
    \item \textbf{Total upvotes received} on songs added by the user
    \item \textbf{Total downvotes received} on songs added by the user
\end{itemize}

Two awards are distributed: \textit{Crowd Favourite} (song with the most upvotes)
and \textit{Top Jammer} (user with the highest reputation delta). Reputation deltas
are computed as:

\begin{equation}
\Delta\text{reputation} = (\text{upvotes\_received} \times 10) -
(\text{downvotes\_received} \times 5) + \text{participation\_bonus}
\label{eq:rep}
\end{equation}

where \texttt{participation\_bonus} is a fixed integer increment awarded for joining
and contributing at least one song. Reputation updates are written to the
\texttt{users} collection in Firestore, persisting the gamification state across
future sessions.

%% ─────────────────────────────────────────────
%%  COMMENTED OUT — sections to be written next
%%  + old synopsis section
%% ─────────────────────────────────────────────


\section{Results and Discussion}\label{sec:results}

\subsection{Functional Test Setup}

Functional evaluation was conducted across two physical Android devices: one running
API level 33 (Android 13) acting as the host, and one running API level 25
(Android 7.1.1) acting as a guest participant. Both devices were authenticated to
separate Firebase user accounts within the same Firebase project and connected to the
same Wi-Fi network. No network throttling was applied. All scenarios were executed
at least twice to confirm reproducibility.

\subsection{Test Scenarios and Outcomes}

Table~\ref{tab:results} summarises the functional test scenarios executed and their
observed outcomes.

\begin{table*}[!ht]
\centering
\caption{Functional test scenarios and observed outcomes}
\label{tab:results}
\begin{tabular}{p{3.0cm} p{5.2cm} p{4.6cm}}
\toprule
\textbf{Scenario} & \textbf{Description} & \textbf{Observed Outcome} \\
\midrule

Room creation and joining
& Host creates a room and receives a unique alphanumeric room code. Guest enters
  the code on a second device.
& Both devices entered \texttt{QueueActivity} with consistent room name and code.
  Active user count incremented on both screens within one second. \\

\addlinespace

Song search and addition
& Guest searches for a song by title and taps a result to add it.
& YouTube Data API returned results within two to three seconds. Song appeared in
  the queue on both devices within one to two seconds of addition. \\

\addlinespace

Vote-based queue reordering
& Two songs present in queue. Guest upvotes the lower-ranked song.
& Song score incremented live on both devices. The upvoted song moved to the top
  of the queue on both devices without a manual refresh. \\

\addlinespace

Host playback and skip
& Host presses play. After several seconds, host presses skip.
& Audio began on the host device. On skip, the next highest-scored song was
  selected and playback began from its start on both devices. \\

\addlinespace

Non-host playback synchronisation
& Host begins playback. Guest device is observed.
& Guest device began playback of the same song within one to two seconds,
  synchronised via the \texttt{nowPlaying} Firestore field. \\

\addlinespace

SFW content filtering
& Host enables SFW mode. Guest attempts to add a song flagged
  \texttt{ytAgeRestricted}.
& Song appeared in search results. On tap, a Snackbar message was shown and
  the song was not added to the queue. No Firestore write was made. \\

\addlinespace

History view and re-queue
& After songs have played, guest navigates to the history tab and taps a song.
& History listed played songs in reverse chronological order. Tapped song was
  re-added with a reset score and attributed to the guest. Queue view restored
  automatically. \\

\addlinespace

Double vote prevention
& A user attempts to vote on the same song a second time.
& The Firestore transaction detected the existing vote document and aborted.
  Song score was unchanged. \\

\addlinespace

Presence and disconnect
& Guest closes the application without explicitly leaving the room.
& RTDB \texttt{onDisconnect()} removed the guest's presence entry. The active
  user count on the host's screen decremented within the RTDB disconnect
  detection window. \\

\addlinespace

Session end and reputation
& Host ends the session from room settings.
& Both devices navigated to \texttt{SessionSummaryActivity}, displaying the
  crowd favourite and top jammer. Reputation deltas were written atomically to
  the \texttt{users} collection. \\

\bottomrule
\end{tabular}
\end{table*}

\subsection{Discussion}

\subsubsection{Real-Time Synchronisation}

Across all multi-device test scenarios, Firestore snapshot listeners propagated
state changes — song additions, vote updates, queue reordering, and now-playing
transitions — to all connected clients within one to two seconds on Wi-Fi. This is
consistent with the expected behaviour described in Section~\ref{sec:methodology}
and confirms that the serverless Firestore architecture is sufficient for real-time
collaborative queue management at small room sizes. No missed updates or stale
states were observed under two concurrent clients.

\subsubsection{Vote-Based Queue Ordering}

The score-based ordering mechanism functioned correctly in all test cases. Upvoting
a song incremented its score and the Firestore \texttt{orderBy("score", DESCENDING)}
query immediately delivered an updated, correctly ordered result set to both devices.
When the host pressed skip, \texttt{reorderAndPopNextSong()} selected the
highest-scored song, confirming that the front-end ordering and the back-end song
selection are consistent. Songs with equal scores were ordered by \texttt{addedAt}
ascending, which was verified by adding two songs in sequence at score zero and
confirming the earlier addition played first.

\subsubsection{Playback Synchronisation}

Host-controlled playback via the YouTube IFrame API operated correctly for standard
videos. Non-host playback synchronisation, driven by the \texttt{nowPlaying} field,
began within the Firestore propagation window. A confirmed limitation is that
age-restricted YouTube videos fail silently with player error codes 101 or 150 and
are skipped automatically by the error handler, as described in
Section~\ref{sec:limitations}.

\subsubsection{SFW Content Filtering}

The SFW filter correctly blocked \texttt{ytAgeRestricted} songs from being added when
active. Rejection was instantaneous and client-side, with no Firestore write
attempted. Songs still appeared in search results in SFW mode, allowing users to
discover them, but addition was blocked at the point of user intent. As noted in
Section~\ref{sec:limitations}, the filter is best-effort and dependent on YouTube's
content labelling coverage.

\subsubsection{Presence and Session Management}

The RTDB presence mechanism correctly reflected the active user count across all
tested transitions: joining, leaving via the back button, and disconnecting by closing
the application. The \texttt{onDisconnect()} handler demonstrated correct behaviour
on unclean disconnection, though removal latency varied between approximately five
and thirty seconds, consistent with RTDB's server-side disconnect detection window.

\subsubsection{Gamification}

Session-end reputation computation executed correctly in all tests. The crowd
favourite and top jammer were correctly identified and displayed in
\texttt{SessionSummaryActivity}. Reputation increments were verified to be written
atomically to the \texttt{users} collection alongside the room status update in a
single Firestore batch, preventing partial state.

\clearpage
\section{Conclusion}\label{sec:conclusion}

This paper presented JamVote, a native Android application implementing a
non-blocking democratic music queue system built on Firebase Firestore and Firebase
Realtime Database without a dedicated application server. The system was designed to
address the following research question: \textit{How can a mobile application
implement a non-blocking, democratic music queue system that dynamically balances
collective engagement, user fairness, and playback continuity without requiring
unanimous participation or a centralised moderator?}

The implemented system demonstrates that this is achievable within a lightweight,
serverless mobile architecture. Songs are added to a shared queue and ranked in real
time by a vote score derived from upvotes and downvotes, with Firestore's
\texttt{orderBy} mechanism propagating queue reordering to all connected clients on
every vote. Playback continuity is guaranteed by the host client, which automatically
advances the queue when a song ends and can skip the current song at any time.
Non-host clients synchronise playback by observing the \texttt{nowPlaying} field on
the room document, achieving consistent audio state across devices within Firestore's
propagation latency of approximately one to two seconds. Vote integrity is enforced
at the database level through a composite-key deduplication scheme that prevents
double voting without requiring a query. A content moderation layer blocks
age-restricted content from being added to the queue when SFW mode is active. A
session-end gamification module computes per-user reputation deltas and identifies
session awards, persisting them to a cross-session user profile in Firestore.

Functional evaluation across two physical Android devices confirmed correct behaviour
across all principal features: room creation and joining, song search and addition,
vote-based queue reordering, host and non-host playback synchronisation, SFW content
filtering, history viewing and re-queuing, double vote prevention, presence tracking,
and session-end reputation computation. The principal confirmed limitation is the
inability to play age-restricted YouTube videos in the embedded \texttt{WebView}
without an authenticated YouTube browser session, which causes such videos to be
silently skipped.

The results confirm that the combination of Firestore's real-time snapshot listener
infrastructure and RTDB's \texttt{onDisconnect()} mechanism is sufficient to sustain
democratic, presence-aware queue governance across concurrent mobile clients without
server-side business logic, at the scale of small private listening rooms.

\section{Study Limitations}\label{sec:limitations}

\subsection{Functional Testing Scope}

The evaluation of JamVote was conducted through functional testing across less than 5 physical Android devices and did not include systematic load testing or stress
testing under high concurrency. The behaviour of the system under large room sizes, specifically queue reordering and presence heartbeat
aggregation, was not evaluated beyond three simultaneous clients. It is therefore not possible to claim that the current implementation scales to rooms with tens or
hundreds of concurrent participants without further testing. Firestore read and write
quotas under the free Spark tier impose hard limits that were not approached during
testing but would become relevant in production deployments.

\subsection{Device and Platform Constraints}

Testing was conducted on three devices running API levels 25 and 33 respectively,
which does not constitute a representative sample of the Android device ecosystem.
Manufacturer-specific battery optimisation policies on certain Android devices
aggressively kill background processes, which may cause the presence heartbeat
handler to be suspended, leading to stale active user counts. This behaviour was not systematically evaluated across device
manufacturers. iOS support was outside the scope of this study; the application is Android-exclusive and no cross-platform evaluation was performed.

\subsection{Network Environment Restrictions}

All functional testing was conducted on a standard Wi-Fi network with no throttling, packet loss simulation, or latency injection. The real-time synchronisation latency of one to two seconds reported for Firestore snapshot
listeners reflects idealised network conditions. Performance under high-latency
or intermittent connectivity conditions common in party or event environments
where cellular networks may be congested was not evaluated. The
\texttt{onDisconnect()} handler that closes the room on host disconnection
relies on Firestore's server-side detection of client dropout, which may fire
with significant delay under degraded network conditions, leaving guest clients
in an inconsistent state.

\subsection{Content Filtering Accuracy}

The SFW mode implementation relies on YouTube's \texttt{ytAgeRestricted} content
rating flag, which is acknowledged as a best-effort filter. YouTube does not
systematically label all explicit music content, and a substantial proportion of
explicit tracks do not carry the age-restriction flag. This means the SFW mode
cannot provide a reliable content guarantee and should be considered an
approximate filter rather than an enforceable policy. A production-grade
implementation would require integration with a dedicated content moderation
API or the Spotify \texttt{track.explicit} field, which provides more reliable
labelling.

A further limitation affects age-restricted videos outside of SFW mode. YouTube's
embedded player requires an authenticated browser session to play
\texttt{ytAgeRestricted} content. Because the application's \texttt{WebView} maintains
no YouTube user session, such videos fail silently with player error codes 101 or 150
and are skipped by the error handler. This cannot be resolved without establishing a
YouTube browser session within the application, which is outside the current
implementation scope.

\subsection{Host Trust and Security Model}

The current architecture places significant trust in the host client. Queue
reordering and session-end reputation computation are both executed client-side on
the host device and written directly to Firestore. A malicious or modified host
client could manipulate queue order or issue incorrect reputation deltas without
detection.
Moving these operations to Firebase Cloud Functions would eliminate this trust
assumption but was outside the scope of the current implementation. The Firestore
security rules deployed for this study enforce per-user write restrictions on
votes and presence documents but do not validate the correctness of sort or
scoring operations.

\subsection{Gamification Evaluation}

The gamification and reputation system was implemented and verified for functional
correctness, badges are awarded and reputation deltas are written correctly at
session end but was not evaluated for its effect on user behaviour or engagement.
Claims made in the related literature regarding the positive effect of reputation
systems on sustained engagement~\cite{jamvote_online_community_motivation,
jamvote_gamification_works} were not empirically validated within this app.
A user study measuring curation quality, session duration, or return participation
across sessions with and without the gamification layer was outside the project
scope.


\section{Future Scope}\label{sec:future}

Several directions for extending and improving JamVote are identified based on the
scope boundaries and design decisions of the current implementation.

\subsection{Methodological Improvements}

The current dynamic scoring formula assigns a fixed weight to contributor reputation
as a tiebreaker. Future work could explore adaptive weighting mechanisms that adjust
the reputation coefficient based on session context, such as room size or historical
voting variance, potentially improving queue fairness in large heterogeneous groups.

\subsection{Spotify Integration and Music Provider Abstraction}

The current implementation uses the YouTube Data API v3 for music search and playback.
Spotify was not adopted as the music backend because its Web API requires an extensive
OAuth permission scope, mandates application review and approval for production access,
and its playback SDK is a paid service restricted to Spotify Premium subscribers,
making it impractical for an open, freely accessible application at this stage.

A natural extension of the current implementation would nonetheless be to introduce a
music provider abstraction layer and swap the existing \texttt{YouTubeProvider} for a
Spotify-backed implementation. Spotify integration would provide reliable explicit
content flagging via the \texttt{track.explicit} boolean field, gapless playback, and
access to audio features such as tempo, energy, and danceability, which could be used
to implement a vibe continuity engine that automatically suggests songs matching the
current session's acoustic profile.

\subsection{Vibe Continuity and Automated Queue Filling}

When the queue is exhausted, the current implementation stalls playback. A vibe
continuity module could automatically suggest or insert songs based on the acoustic
similarity of recently played tracks, using audio feature vectors from the Spotify
Web API or a lightweight on-device genre classifier. This would ensure uninterrupted
playback in sessions where user participation drops temporarily.

\subsection{Distributed Host and Session Resilience}

The current architecture assigns permanent host authority to the room creator, and
closes the session on host disconnection. A distributed host election mechanism,
for example, assigning host authority to the participant with the lowest join
timestamp among active users would improve session resilience in long-running or
large-group settings where host dropout is more probable.

\subsection{Scalability and Public Room Support}

The current implementation targets private rooms with small participant counts.
Scaling to public rooms with potentially hundreds of concurrent users would require
rate-limiting vote writes, batching presence heartbeats, and potentially migrating
queue state management to Firebase Cloud Functions to reduce the trust placed in
the host client for sort and dominance operations. This would also remove the
current vulnerability in which a malicious host client could manipulate queue
ordering.

\subsection{Manipulation Resistance}

As noted in the gamification literature~\cite{jamvote_online_community_motivation,
jamvote_gamification_works}, reputation systems in competitive voting environments
are susceptible to coordinated manipulation. Future work could investigate
rate-limiting vote operations per user per session, detecting anomalous voting
patterns via Firestore Cloud Functions, and applying trust-decay mechanisms that
reduce the influence of reputation scores accumulated through suspected collusion.

\subsection{Cross-Platform Extension}

The current application targets Android exclusively. A cross-platform implementation
using a framework such as Flutter would extend JamVote to iOS, enabling mixed-device
sessions and broadening the potential user base for party and event use cases.
Web-based room management, allowing non-mobile participants to join and vote via a
browser, is a natural extension that would not require changes to the core Firestore
data model.

\begin{comment}
\section{Project Synopsis}\label{sec1}

\subsection{Purpose of The App}\label{subssec2}

JamVote is a real-time collaborative music streaming platform designed to enhance shared
listening experiences in parties, remote friend groups, and public community rooms.

Unlike existing collaborative playlist systems, JamVote introduces an intelligent,
non-blocking governance model that balances democratic input, fairness, and automated vibe
continuity, ensuring sessions never stall and no single user dominates the queue.

\subsection{Problem Statement}\label{subsec2}

Current collaborative music solutions allow multiple users to add and reorder songs in a
shared queue. However, these systems exhibit several limitations:
\begin{itemize}
    \item \textbf{Manual Queue Dominance}: Users can drag and reorder tracks arbitrarily,
    leading to unfair control.
    \item \textbf{No Collective Feedback Mechanism}: There is no structured voting or
    engagement-based ranking.
    \item \textbf{No Continuity Intelligence}: When the queue runs out, playback depends
    on an individual's existing playlist.
    \item \textbf{Lack of Gamification}: No reward system for good music curation.
    \item \textbf{No Adaptive Governance}: No differentiation between active and inactive
    users.
\end{itemize}

\subsection{Objectives}\label{subsec2b}

The objectives of JamVote are:
\begin{itemize}
    \item To design and develop a real-time collaborative music Android application.
    \item To design a non-blocking collaborative queue system.
    \item To implement a dynamic scoring mechanism.
    \item To ensure fairness by preventing queue monopolization.
    \item To develop a genre continuity engine for auto-filling music when required.
    \item To introduce gamification and reputation systems.
    \item To maintain scalability across both private and public listening rooms.
\end{itemize}

\subsection{Project Scope}\label{subsec2c}

\subsubsection{User Authentication and Room Management}
\subsubsection{Real-Time Collaborative Queue System}
\subsubsection{Non-Blocking Voting Mechanism}
\subsubsection{Vibe Continuity Module}
\subsubsection{Gamification and Reputation System}

\end{comment}

\clearpage
\bibliography{sn-bibliography}

\end{document}
